% Preface: the README's About, During the Course, Learning Outcomes and Written Papers sections,
% the course plan from info/README.md, and L01's prerequisites, as the front matter of a book.

\chapter{Preface}
\markboth{Preface}{Preface}

This book is for experienced C developers who want to learn how modern C++ can be used
effectively in embedded systems. It covers the development of embedded systems in modern C++
(C++17), with a focus on practical design techniques for resource-constrained systems.

Unlike an introductory C++ text, it assumes strong prior experience with C, and it covers only the
C++ concepts that are useful in embedded systems:
\begin{itemize}
  \item The transition from C to modern embedded C++.
  \item Compile-time programming (\code{constexpr}, templates).
  \item Deterministic embedded design (\code{noexcept}, controlled use of dynamic memory
    allocation).
  \item Structs and classes for hardware drivers.
  \item Encapsulation and object-oriented design in embedded systems.
  \item Interfaces and polymorphism.
  \item Factory patterns and dependency management.
  \item Templates and generic programming.
  \item Copy and move semantics in embedded C++.
  \item Template utilities for low-level programming.
  \item Threading and synchronization primitives.
  \item Designing reusable and testable driver libraries.
\end{itemize}

\textbf{Along the way you will write code, not just read it.} You will implement simple drivers
using modern C++, apply compile-time techniques to low-level programming, refactor traditional C
patterns into safer C++ constructs, and design modular embedded software using interfaces and
factories, always with an eye on how C++ can improve modularity and code structure
\emph{without} introducing runtime overhead. The examples and exercises stay close to typical
embedded problems:
\begin{itemize}
  \item Bit manipulation utilities.
  \item Register access helpers.
  \item Simple driver abstractions.
  \item Compile-time utilities for low-level code.
  \item Interface-based driver design.
  \item Template-based utilities.
\end{itemize}

\section*{What you should be able to do at the end}
Having worked through the book, you should be able to:
\begin{itemize}
  \item Use modern C++ features effectively in embedded systems.
  \item Write low-level code using templates and compile-time techniques.
  \item Design simple hardware drivers using structs and classes.
  \item Use interfaces to decouple system logic from hardware implementations.
  \item Apply factory patterns to construct embedded components.
  \item Use templates to build reusable embedded utilities.
  \item Apply object-oriented design without sacrificing performance or determinism.
  \item Understand how C++ abstractions translate to machine code in embedded environments.
\end{itemize}

\section*{What you are assumed to know}
The subject of this book is \textbf{the C++ that pays its way on a microcontroller}, not embedded
programming in general. You are expected to arrive already comfortable with:
\begin{itemize}
  \item C programming: functions, pointers, structs and the rest of the language.
  \item Bit manipulation, and the embedded programming concepts it is used for.
\end{itemize}
None of that is taught from scratch. What the book does teach from nothing is every C++ feature it
uses, starting in Chapter~1 with the difference between a \code{.h} and a \code{.hpp} file.

\section*{How the book is organised}
The book is typeset from a six-lecture course, and each chapter is one lecture. Its sections are
the lecture's appendices, in the same order and with the same subsections, so a chapter here and a
lecture directory in the repository can be read side by side:

\begin{booktable}{@{}p{5.2em}p{3.6em}L@{}}
  \thead{Chapter} & \thead{Lecture} & \thead{Topic} \\
  \midrule
  1 & \file{L01} & Modern embedded C++ concepts \\
  2 & \file{L02} & Classes \\
  3 & \file{L03} & Inheritance and interfaces \\
  4 & \file{L04} & The factory pattern \\
  5 & \file{L05} & Templates \\
  6 & \file{L06} & Multithreading and synchronization \\
\end{booktable}

Every chapter opens with what is in it, works through the material, closes with a review of what
you should now be able to explain, and ends with its exercises. The appendices hold the two written
papers described below.

\section*{What the exercises look like}
The exercises come in \textbf{sets}. A set builds one small program, or one part of a larger
one, and each exercise in a set adds a piece of it. Every exercise is labelled with its kind:

\begin{booktable}{@{}p{7.4em}L@{}}
  \thead{Kind} & \thead{What it asks of you} \\
  \midrule
  \kindtag{Code} & Write C++, compile it, and check its output against the one the exercise
    gives. \\
  \kindtag[fillaccent2]{Reflection} & Answer in prose: explain what your code does, or why one
    design is better than another. \\
\end{booktable}

\textbf{The solutions are not in this book.} They are published in the companion repository,
under \file{lectures/LNN/appendix/solutions}, grouped the same way as the exercises, each program
with the Makefile that builds and runs it, and with the answers to the reflection questions in the
README beside them. Clone the repository (see \emph{Setting up} below) and you can run a reference
program beside your own. Try each exercise before you read its solution.

\section*{Two written papers, and what they are not}
Nothing in the course this book comes from is marked. Assessment is the exercise set after every
chapter, each with a published solution, and the questions in every chapter's review.

Appendices~\ref{exam:ch:about} to \ref{ansb:ch:answers} hold two three-hour papers with worked
solutions, and they check something else: \textbf{what you can reconstruct on paper, with nothing
in front of you.} Eight questions each, mixing theory with C++ you either read and repair or write
from scratch, and no compiler in the room to tell you which it is.

\textbf{They exist purely so that you can test your own knowledge after working through the book.
They gate nothing, they are not a qualification, and no part of the book requires them.}

\textbf{Take one after the last chapter.} Both papers draw on all six chapters, so sitting one
partway through examines material you have not read yet.

\section*{Setting up}
Everything in this book runs in a Linux terminal: on Linux itself, or on Windows through WSL. You
need three things:
\begin{itemize}
  \item \textbf{GCC and G++}, to compile the example programs.
  \item \textbf{Linux or WSL}, to compile and run them in a terminal environment.
  \item \textbf{An editor}; Visual Studio Code is the one the course uses, available at
    \url{https://code.visualstudio.com/download}.
\end{itemize}
\Secref{c1:sec:compilation} walks through installing WSL, Ubuntu and the compiler, and writing the
Makefile every program in the book is built with. Then clone the companion repository:

\begin{shell}
git clone https://github.com/qrtech-academy/modern-embedded-cpp.git
\end{shell}

\noindent Two commands at its root cover everything the book asks of it:

\begin{shell}
make build               # Build every demo and solution that has a Makefile.
make format-check        # Check the formatting of every C++ file.
\end{shell}
